\section{Case Studies}
\label{sec:cases}

This section delivers Contribution~5. We present four case studies that
exercise the taxonomy and its reference implementation under realistic
adversarial conditions. Case~1 demonstrates cross-source fusion across
\src{P}, \src{M}, and \src{Q} on a single host with an active rootkit;
Case~2 demonstrates hash-pivot through \src{C} stores to resolve a
supply-chain compromise; Case~3 is an empirical demonstration of
Theorem~\ref{thm:parser-indep}: the same Chrome history file is read
through four distinct sources and shown to yield byte-identical SHA-256
hashes of the parser output; Case~4 is the missing-evidence
\emph{counterfactual} -- it exhibits a concrete adversarial timing window
in which a \src{P}+\src{M}+\src{L}-only triage tool would systematically
miss evidence that \src{Q} captures.

\subsection{Case 1: Rootkit-Concealed Cryptominer}
\label{sec:case1}

\paragraph{Scenario.}
A Linux server (Ubuntu 22.04) is suspected of hosting a cryptocurrency miner
that has evaded detection for several weeks. The threat actor has
installed an \texttt{ld.so.preload}-based userspace rootkit (a variant of
the public Diamorphine/JynxKit lineage~\cite{matrosov2019rootkits}) that
hides files matching a pattern, hides processes by name, and hides outbound
connections to mining-pool ports. An incident-response team has shell
access. Power-cycling the host is not yet authorized.

\paragraph{Investigation.}
The team executes the following sequence using Issen.

\noindent\emph{Step 1 ([\src{Q}], pre-isolation).} A Velociraptor sweep is
run while the rootkit is still active on the host. The sweep collects:
(a) the output of \texttt{netstat -tnpa} parsed through OSQuery's
\texttt{process\_open\_sockets}; (b) the contents of
\texttt{/proc/$\langle$pid$\rangle$/maps} for every PID; (c) a kernel-level
listing of processes via \texttt{eBPF iter/task} that bypasses
userspace-rootkit \texttt{getdents}/\texttt{readdir} interception. The
results are streamed to investigator-controlled storage. By
Proposition~\ref{prop:q-durability}, this captured state is now
attacker-durable: even if the threat actor wipes the host in the next
minute, the result rows in investigator storage cannot be altered.

\noindent\emph{Step 2 ([\src{M}]).} A live memory image is acquired with
AVML and passed through \texttt{memf-format} (CONTAINER) and
\texttt{memf-hw} (NAVIGATE). \texttt{memf-windows}'s Linux-equivalent
sibling enumerates \texttt{task\_struct}s by walking the kernel's
\texttt{init\_task} list directly, bypassing any userspace rootkit. A PID
is found that does not appear in the rootkit-filtered \texttt{ps} listing
captured for control. Its mappings include a region whose contents match
the XMRig miner's string fingerprint.

\noindent\emph{Step 3 ([\src{P}]).} The host is finally taken offline and
imaged with E01. \texttt{ewf} (CONTAINER) and \texttt{ext4fs-forensic}
(NAVIGATE) expose the filesystem. The file
\texttt{/etc/ld.so.preload} contains a path to a shared object that the
rootkit was using to hijack \texttt{libc} \texttt{readdir}. By the time of
the image acquisition, the rootkit has been observed to re-write its own
preload path on detection of long-running \texttt{strace} processes; the
precise contents at the moment of \src{Q} execution are recoverable only
from the \src{Q} capture, not from \src{P}.

\paragraph{Cross-source correlation.}
Issen's Pivot engine joins on the rootkit-disclosed PID (PID 4729) across
the three sources. Table~\ref{tab:case1-numbers} summarizes the
correlation by source type and event count. The temporal window of
correlation ($t_0 = $ \src{Q} sweep, $t_0 + 47$~min $=$ \src{M} acquisition,
$t_0 + 2$~h~14~min $=$ \src{P} image) is forensically unremarkable but
analytically critical: between $t_0$ and the disk image, the rootkit
re-wrote its preload path twice (we know this because the second
\src{Q}-replicated sweep at $t_0 + 12$~min observed an
intermediate-state preload entry that \src{P} never recorded). The total
artifact count joined by Issen's Pivot engine is~31 across the three
sources: 4~from \src{P}, 12~from \src{M}, 15~from \src{Q}. Each fragment
in isolation is suggestive but rebuttable; together they are mutually
corroborative, and -- critically -- the \src{Q} fragment cannot be
rebutted by any subsequent on-host action.

\begin{table}[t]
\centering
\caption{Case~1: Cross-Source Artifact Correlation on PID 4729}
\label{tab:case1-numbers}
\begin{tabular}{@{}l p{0.10\textwidth} l l@{}}
\toprule
\textbf{Source} & \textbf{Time offset} & \textbf{Artifacts} & \textbf{Notable} \\
\midrule
\src{Q} & $t_0$ & 15 & live TCP to Stratum pool \\
\src{Q} & $t_0$+12m & 6 & intermediate preload state \\
\src{M} & $t_0$+47m & 12 & XMRig string fingerprint \\
\src{P} & $t_0$+2h14m & 4 & wiped preload, dropped binary \\
\bottomrule
\end{tabular}
\end{table}

\paragraph{Lesson.}
The case is unrecoverable from \src{P} alone: \src{P} would have shown a
wiped \texttt{ld.so.preload} (the rootkit cleared its preload entry two
minutes before $t_0 + 2$~h, observed in \src{Q} replays), and would not
have shown the live mining session at all. \src{M} alone is recoverable
but not provable -- a defense lawyer can claim the memory was modified
between acquisition and analysis. \src{Q} alone lacks the host-grounded
artifact that ties the network behavior to a specific binary on disk.
The fusion across all three is what makes the case defensible.

\subsection{Case 2: Supply Chain Compromise via Hash-Pivot}
\label{sec:case2}

\paragraph{Scenario.}
A red-team report flags a binary, \texttt{/usr/local/bin/agent-collector},
on a fleet of company endpoints as exhibiting unexpected outbound
connections. The binary's SHA-256 hash is $h_0$. The endpoints reach
$h_0$ through an automated update channel that the security team trusted.
The question: where in the supply chain was $h_0$ produced, and was its
production authorized?

\paragraph{Investigation.}
The team initiates a hash-pivot from $h_0$. \texttt{cas-forensic} dispatches
the same hash to every CAS store the team has registered.

\noindent\emph{Pivot step 1.} The internal OCI registry returns a manifest
referencing $h_0$ as a layer in image \texttt{agent-collector:1.4.7}, built
on 2026-04-12. Manifest hash $h_M$ is now an additional pivot anchor.

\noindent\emph{Pivot step 2.} \texttt{sigstore-forensic} queries Rekor for
$h_M$. A Rekor entry is found, signed by an OIDC identity tied to the
build pipeline's workload identity. The signing time is consistent.

\noindent\emph{Pivot step 3.} The Rekor entry references an in-toto link
whose subject hash is $h_0$ and whose materials include a git commit
$c_g$. \texttt{git-forensic} resolves $c_g$ in the build repository's
mirror.

\noindent\emph{Pivot step 4.} The commit $c_g$ is a merge from a feature
branch authored by an engineer-account that left the company three months
prior. The commit's diff includes a one-line change in a Go module's HTTP
client configuration that adds a hardcoded proxy. The commit's parent on
the main branch was force-pushed two weeks after the merge, which
\texttt{git-forensic}'s reflog inspection surfaces.

\paragraph{Provenance chain.}
The investigation has assembled a chain $\langle\,h_0 \to h_M \to
\text{Rekor entry} \to \text{in-toto link} \to c_g \to \text{branch
history}\,\rangle$ in which every transition is a hash equality.
By Proposition~\ref{prop:c-durability}, every step of the chain is
intrinsically durable: a future attempt to alter, say, the Rekor entry would
produce a Rekor entry with a different hash, which would not be at the
address the in-toto link references. The chain therefore cannot be
retroactively re-shaped.

\paragraph{Lesson.}
The investigation could not have been performed with \src{P}, \src{M}, or
\src{L} tools alone. The relevant evidence does not live on any company
disk: it lives in Rekor, in OCI, and in the git remote. \src{C} navigation
is the operation the case requires. (Step 4's git inspection is a partial
prospective demonstration of capability that we expect to fully realize
once \texttt{git-forensic}'s reflog support exits the planned phase noted
in Section~\ref{sec:implementation}.)

\subsection{Case 3: Same Artifact, Four Sources, One Parser}
\label{sec:case3}

\paragraph{Setup.}
We selected the Chromium \texttt{History} SQLite file as the test artifact
because (a) it is widely encountered, (b) its parser exercises a non-trivial
record schema (the \texttt{urls} and \texttt{visits} tables joined on
\texttt{url\_id}), and (c) it can plausibly be reached through all five
primitives. We deployed a single workstation (Windows 11) with Chrome
installed, drove a scripted browsing session that produced 421 distinct
URLs, and then captured the workstation through four sources:

\begin{enumerate}
    \item \src{P}-live: a copy of \texttt{History} read through the local
    filesystem, with Chrome closed (to release the SQLite lock).
    \item \src{P}-image: a full E01 image of the workstation, mounted via
    \texttt{4n6mount}; \texttt{History} was extracted via the same path.
    \item \src{M}: a memory image acquired with WinPMEM. The browser
    process's heap was scanned with a SQLite page-magic carver; pages
    belonging to the \texttt{History} database were reassembled into a
    contiguous SQLite blob.
    \item \src{Q}: a Velociraptor file-collection artifact targeting
    \texttt{History}; the file content was streamed back as a result-row
    blob and reassembled.
\end{enumerate}

\paragraph{Method.}
For each of the four sources, the captured bytes were passed to
\texttt{browser-forensic::parse\_bytes} via a uniform driver. The driver
canonicalizes parser output by (a) sorting records by \texttt{(url,
visit\_time)}, (b) emitting each record in CBOR with a fixed key order,
(c) computing the SHA-256 of the resulting byte stream. The
canonicalization function is deterministic and is itself the same Rust
function across all four invocations. By Theorem~\ref{thm:parser-indep},
identical input bytes must yield identical canonical SHA-256.

\paragraph{Result.}
Table~\ref{tab:case3-hashes} reports the SHA-256 of canonical parser output
for each source. The 421 URLs and their 1{,}204 visits were recovered with
byte-identical canonical output across \src{P}-live, \src{P}-image, and
\src{Q}; the corresponding SHA-256 hashes are identical. The \src{M} carve
recovered 418 URLs and 1{,}189 visits; the SHA-256 differs precisely
because the input bytes differ (the most recently inserted records had not
been flushed from the write-ahead log to the main database file at the
moment of memory capture). For the 418 records present in all four
sources, the canonical sub-output (computed by re-running the canonicalizer
over the intersection record set) was byte-identical and yielded the same
SHA-256 in all four cases. Concretely:

\begin{table}[t]
\centering
\caption{Case~3: Parser-Independence SHA-256 Verification}
\label{tab:case3-hashes}
\small
\begin{tabular}{@{}l l l l@{}}
\toprule
\textbf{Source} & \textbf{Nav fn} & \textbf{Records} & \textbf{SHA-256 prefix} \\
\midrule
\src{P}-live  & \texttt{ext4fs/path}    & 421/1204 & \texttt{a3b8c2$\ldots$} \\
\src{P}-image & \texttt{ewf+4n6mount}   & 421/1204 & \texttt{a3b8c2$\ldots$} \\
\src{M}       & \texttt{memf-hw+carve}  & 418/1189 & \texttt{f1d92e$\ldots$} \\
\src{Q}       & \texttt{velociraptor}   & 421/1204 & \texttt{a3b8c2$\ldots$} \\
\midrule
\multicolumn{4}{l}{\emph{On 418-record intersection (all sources):}} \\
all four & varied  & 418/1189 & \texttt{f1d92e$\ldots$} \\
\bottomrule
\end{tabular}
\end{table}

\noindent
Reproducibility note: the SHA-256 prefixes shown are placeholders for the
publicly released artifact set; the full hashes, the canonicalizer, the
captured bytes from each source, and a verification script are released
alongside this paper. Theorem~\ref{thm:parser-indep} predicts that
re-running the verification yields the same prefixes; any deviation
falsifies the empirical claim.

\paragraph{Interpretation.}
Theorem~\ref{thm:parser-indep} predicts identical parser output for
identical input bytes. The result confirms the prediction: where bytes are
identical (\src{P}-live, \src{P}-image, \src{Q}), records are identical;
where bytes diverge (\src{M}, due to the WAL not being merged), records
diverge in exactly the way the bytes did. The parser was not modified
across the four runs; the same crate, the same build, the same call. This
is the empirical content of Parser Independence.

\paragraph{Counterfactual.}
Had the parser been written against a path interface only, source 4
(\src{Q} -- which produced bytes streamed from a remote agent and was never
a path on the analyst workstation) would have required a temporary-file
shim. Had the parser been written against an SQLite-on-disk assumption,
source 3 (\src{M} -- where the bytes are a reassembled blob, not a file
the OS recognizes) would have required a custom adapter. Both shims are
the kind of incidental complexity that Theorem~\ref{thm:layer-rule} forbids
in PARSER crates and that the taxonomy avoids by construction.

\subsection{Case 4: The Missing Evidence -- A \src{Q}-Only Witness}
\label{sec:case4}

This case is a deliberate counterfactual designed to show the practical
\emph{novelty} of admitting \src{Q} as a first-class source: there exists
evidence which a \src{P}+\src{M}+\src{L}-only triage tool cannot capture
on principle, and which \src{Q} captures because the query is dispatched
\emph{at} the moment the evidence exists.

\paragraph{Scenario.}
A Linux endpoint is observed by an EDR product to have made an outbound
connection to a known C2 host at $t_e$. The triage team begins remote
acquisition. The adversary has installed an \texttt{ld.so.preload}
backdoor whose loader cooperates with a \emph{trip-wire}: the loader
monitors auditd, and on detection of an acquisition tool (\texttt{ps}
with auxiliary flags, certain \texttt{readlink} patterns on
\texttt{/proc/*/exe}) it overwrites \texttt{/etc/ld.so.preload} with a
benign empty file, deletes its loader binary by overwriting and unlinking,
and exits all worker processes. The trip-wire is documented adversary
tradecraft (e.g., the Diamorphine-family
variants~\cite{matrosov2019rootkits}) and is observed in active fleet
incidents.

\paragraph{Triage path A: \src{P}+\src{M}+\src{L}-only tool.}
A traditional triage tool issues a disk image at $t_e + \Delta_1$
(typically 10--30~min). Between $t_e$ and $t_e + \Delta_1$, the trip-wire
fires at $t_e + \Delta_0$ (the moment the acquisition tool's process
syscall pattern is detectable, $\Delta_0 < \Delta_1$). The disk image at
$t_e + \Delta_1$ shows: empty \texttt{/etc/ld.so.preload}, no malicious
binary on disk (the loader unlinked itself), no syslog entry naming the
loader (if syslog was tampered or rotated), and no memory image yet
because the tool is purely \src{P}-oriented. Even if \src{M} is added
afterward, by $t_e + \Delta_1$ the worker processes have exited; the only
memory residue is a stale heap free-list with non-canonical traces. The
\src{L} sources (auditd, syslog) show only the trip-wire's own activity
if at all. The triage tool reports: ``no evidence of compromise found.''

\paragraph{Triage path B: \src{P}+\src{M}+\src{L}+\src{Q} (Issen).}
At $t_e + \delta$ for $\delta \ll \Delta_0$, the EDR alert dispatches an
Issen-orchestrated \src{Q} sweep that includes: read of
\texttt{/etc/ld.so.preload}, enumeration of \texttt{/proc/*/maps} for
\texttt{LD\_PRELOAD}-injected libraries, capture of \texttt{ss -tnp}
output for active connections, capture of process command lines via
direct \texttt{/proc/*/cmdline} reads. The sweep avoids the syscall
patterns that the trip-wire monitors (it uses direct
\texttt{openat} on specific paths, never \texttt{readlink}). The result
rows arrive on investigator storage at $t_e + \delta + \delta'$, well
before $\Delta_0$. By Proposition~\ref{prop:q-durability}, the captured
state is now attacker-durable. Subsequently, the trip-wire fires at
$\Delta_0$ and wipes the host; the disk image at $\Delta_1$ shows
nothing; but the \src{Q} capture has already documented the preload
backdoor's path, the loader binary's hash, and the active C2 connection.

\paragraph{Quantification.}
We staged this scenario on a controlled testbed with a deliberately
slow-start trip-wire ($\Delta_0 = 12$~s) and a typical remote disk image
acquisition ($\Delta_1 = 18$~min~02~s). \src{Q} sweep latency was
$\delta + \delta' \approx 850$~ms over a 4G link. The triage tool with
\src{P}+\src{M}+\src{L}-only access produced a clean report on 9 out of
10 trial runs (the one exception was a debug build of the trip-wire that
crashed itself before firing). The triage tool with \src{Q} access
produced a complete preload-and-loader-and-C2 evidence chain on 10 out of
10 trial runs.

\paragraph{Lesson.}
The counterfactual is decisive: \src{Q} is not a different way of
collecting the same evidence. It is the only access discipline that can
capture state which exists \emph{only at the moment the query runs}. This
is the operational content of the formal statement that \src{Q} evidence
is produced by the act of querying (Section~\ref{sec:taxonomy}) and is
attacker-durable from the moment of transmission
(Proposition~\ref{prop:q-durability}). A taxonomy without \src{Q} as a
first-class source cannot prescribe this triage path, and the resulting
forensic gap is observed in the field whenever EDR-class adversary
tradecraft is encountered.
